\documentclass[paper=a4,article]{jlreq}

\usepackage[header={time\_log\_app 解説ノート}]{kk}
\usepackage{kk-badge}
\usepackage{listings}
\tcbuselibrary{listings}
\usepackage{tabularray}
\usetikzlibrary{positioning,arrows.meta,calc}

% ---------- 重要度の色 ----------
\definecolor{impA}{HTML}{C0392B}   % 最重要（赤）
\definecolor{impAbg}{HTML}{FBE9E7}
\definecolor{impB}{HTML}{D9822B}   % 重要（橙）
\definecolor{impBbg}{HTML}{FDF1E4}
\definecolor{impC}{HTML}{6F7B85}   % 参考（灰）
\definecolor{impCbg}{HTML}{EEF1F3}
\definecolor{impD}{HTML}{A9B0B6}   % 対象外（薄い灰）
\definecolor{impDbg}{HTML}{F7F8F9}
% ---------- 根拠ラベルの色 ----------
\definecolor{evlog}{HTML}{2F5B82}
\definecolor{evcode}{HTML}{2E7D4F}
\definecolor{evguess}{HTML}{8E44AD}
\definecolor{codebg}{HTML}{F8F9FB}

% 重要度タグ（本文用）
\newtcbox{\kkTag}[1]{on line, colback=#1, colframe=#1, colupper=white,
  arc=1.5pt, boxrule=0pt, boxsep=0pt, left=3pt, right=3pt, top=1.5pt, bottom=1.5pt,
  fontupper=\footnotesize\gtfamily}
\newcommand{\tagA}{\kkTag{impA}{最重要}}
\newcommand{\tagB}{\kkTag{impB}{重要}}
\newcommand{\tagC}{\kkTag{impC}{参考}}
\newcommand{\tagD}{\kkTag{impD}{対象外}}
% 重要度（表用）
\newcommand{\lvA}{\textcolor{impA}{\textgt{■最重要}}}
\newcommand{\lvB}{\textcolor{impB}{\textgt{■重要}}}
\newcommand{\lvC}{\textcolor{impC}{\textgt{■参考}}}
\newcommand{\lvD}{\textcolor{impD}{\textgt{■対象外}}}

% 根拠ラベル
\newtcbox{\kkEv}[1]{on line, colback=white, colframe=#1, colupper=#1,
  arc=1.5pt, boxrule=0.5pt, boxsep=0pt, left=2pt, right=2pt, top=1pt, bottom=1pt,
  fontupper=\scriptsize\gtfamily}
\newcommand{\evLog}{\kkEv{evlog}{根拠：開発ログ}}
\newcommand{\evCode}{\kkEv{evcode}{根拠：コード}}
\newcommand{\evGuess}{\kkEv{evguess}{推測}}

% Python で言うと（注記枠）
\newenvironment{pyaside}{\begin{KKfbox9ao}[Pythonでいうと][1pt][evcode]}{\end{KKfbox9ao}}

% 図の題
\newcommand{\figcap}[1]{\par\smallskip{\centering\small\textgt{#1}\par}\medskip}

% ---------- コードの見た目 ----------
\lstdefinelanguage{TSX}{
  morekeywords={import,from,export,default,function,const,let,return,if,else,type,
    interface,async,await,new,true,false,null,undefined,try,catch,finally,typeof},
  sensitive=true,
  morecomment=[l]{//},
  morecomment=[s]{/*}{*/},
  morestring=[b]',
  morestring=[b]",
  morestring=[b]`,
}
\lstdefinelanguage{plain}{morestring=[b]",morecomment=[l]{\#}}
\lstdefinestyle{kkts}{
  language=TSX,
  basicstyle=\ttfamily\footnotesize,
  keywordstyle=\color{evlog}\bfseries,
  commentstyle=\color{evcode},
  stringstyle=\color{impA!75!black},
  numbers=left, numberstyle=\tiny\color{impC}, numbersep=6pt,
  xleftmargin=1.8em,
  breaklines=true, breakatwhitespace=false,
  postbreak=\mbox{\textcolor{impC}{$\hookrightarrow$}\space},
  columns=fullflexible, keepspaces=true, showstringspaces=false,
  upquote=true, tabsize=2,
}
% 元のコード枠：[追加オプション]{重要度の色}{ファイル}{開始行}
\newtcblisting{codebox}[4][]{%
  enhanced, breakable, listing only,
  colback=codebg, colframe=#2, boxrule=0.5pt, leftrule=2.2mm, arc=1mm,
  left=1mm, right=1mm, top=0.5mm, bottom=0.5mm,
  title={\texttt{#3}\hfill{\scriptsize 元のコード／#4 行目から}},
  fonttitle=\small\gtfamily, coltitle=white, colbacktitle=#2,
  listing options={style=kkts, firstnumber=#4, #1},
  before skip=6pt, after skip=8pt}

% hyperref は必ず最後に読み込む。
\usepackage{hyperref}
\hypersetup{
  luatex,
  pdfencoding=auto,
  colorlinks=true,
  linkcolor=black,
  urlcolor=blue
}

\title{学習記録アプリの中身を読む}
\date{\today}

\begin{document}

\kkcoverspine{TIME LOG APP}{学習記録アプリの中身を読む}{技術・フォルダ・データ・処理の流れを、元のコードで追う}{\today}

\tableofcontents
\newpage

% =====================================================================
\section{DATA}
\begin{description}
  \item[タイトル] 学習記録アプリの中身を読む
  \item[日付] \today
  \item[科目] 自分のアプリ \texttt{time\_log\_app} のコードを理解する
  \item[出典範囲] \texttt{mobile/} のコード（コミット \texttt{5c5039a} の時点）と、\texttt{log/} の開発ログ V1〜V3。
    コードはすべて実際のファイルから行番号つきでそのまま抜き出している。
\end{description}

% =====================================================================
\section{導入}

\subsection{このノートを読み終えたらできること}
\begin{enumerate}
  \item このアプリが「どの技術で」「なぜそれで」作られているかを、自分の言葉で説明できる。
  \item フォルダを見て、どこに何があって、どこから読めばいいかが分かる。
  \item 「タイマーを止めると、ホームの数字が増える」までの流れを、ファイル名と行番号でたどれる。
\end{enumerate}

\subsection{どんなアプリか（おさらい）}
学習時間を記録するアプリ。ただし目的は「もっと勉強させること」ではなく、
\textgt{頑張った分を認めて、疲れているときは休んでいいと伝えること}にある。
開発ログ V1 では「ストリークやランキングのような外からのごほうびをなくし、努力をほめて、サボった日も責めない」、
V3 では「疲労を見えるようにして、あえて休むことを肯定する」と書かれている。\evLog

画面は下のタブ5つ（ホーム・タイマー・カレンダー・疲労・設定）と、ログイン画面でできている。

\subsection{このノートの読み方}

\textgt{色は「どれだけ大事か」を表す。}コードの枠、表、図はすべてこの4色で塗り分けている。

\begin{center}
\begin{tblr}{
  colspec={Q[c,wd=2.0cm] X[l]},
  width=0.92\linewidth,
  hlines={black!15},
  cells={font=\small},
}
\tagA & このアプリ独自の中心部分。\textgt{必ず読む。}ここが分かれば全体がつながる。\\
\tagB & このアプリ独自の部分。中心を読んだ後に、必要になったら読む。\\
\tagC & 雛形（Expo が最初に用意したもの）や設定。今は「ある」と知っていれば十分。\\
\tagD & 自動で作られるもの・秘密の値。\textgt{読まない・触らない。}\\
\end{tblr}
\end{center}

\textgt{「なぜ」には根拠のラベルをつけている。}
\evLog は \texttt{log/} の開発ログに理由が書いてあるもの、
\evCode はコードの中のコメントに理由が書いてあるもの、
\evGuess はどこにも書いていないので、一般的な理由から推測したもの。
推測は、当時の本当の理由と違うことがある。

\textgt{Python と比べる。}開発ログ V1 には「Python でいうと何か」の比較が書かれているので、
このノートでも同じように \textgt{Pythonでいうと} の枠で比べる。

\subsection{全体像}

図1は、このアプリを「層」に分けた地図。上から順に、画面 → 共有データ → 保存場所 と下りていく。
このあとの章は、すべてこの図のどこかの話になる。

\begin{center}
\begin{tikzpicture}[font=\small,
  layer/.style={draw=#1, fill=#1!7, line width=0.9pt, rounded corners=3pt,
    minimum width=8cm, minimum height=1.15cm, align=center, inner sep=3pt},
  side/.style={draw=#1, fill=#1!7, line width=0.9pt, rounded corners=3pt,
    minimum height=1.15cm, align=center, inner sep=3pt},
  arr/.style={-{Stealth[length=2.2mm]}, line width=0.8pt, draw=black!65},
  arr2/.style={{Stealth[length=2.2mm]}-{Stealth[length=2.2mm]}, line width=0.8pt, draw=black!65},
  lab/.style={font=\scriptsize, fill=white, inner sep=1.2pt}]
  \node[draw=black!45, rounded corners=3pt, minimum width=8cm, minimum height=0.8cm, align=center]
    (user) at (0,0) {あなた（スマホのアプリ／ブラウザで開く Web 版）};
  \node[layer=impA] (app) at (0,-1.8)
    {\textgt{画面}\quad\texttt{app/}\\{\footnotesize ホーム・タイマー・カレンダー・疲労・設定・ログイン}};
  \node[layer=impA] (ctx) at (0,-3.9)
    {\textgt{共有データと保存の係}\quad\texttt{context/}\\{\footnotesize 学習時間・前後チェックの記録・疲労・設定・ログイン状態}};
  \node[side=impB, text width=4.0cm] (comp) at (6.7,-1.8)
    {\textgt{画面の部品}\\\texttt{components/}};
  \node[side=impB, text width=4.0cm] (utl) at (6.7,-3.9)
    {\textgt{計算の関数}\\\texttt{utils/}\ \texttt{hooks/}};
  \node[side=impA, text width=2.3cm, minimum height=3.25cm] (typ) at (-6.05,-2.85)
    {\textgt{データの形}\\\texttt{types/}\\[3pt]{\footnotesize 画面も共有データも\\この形に従う}};
  \node[side=impB, text width=3.6cm] (as) at (-2.1,-6.0)
    {\textgt{端末の中}\\{\footnotesize AsyncStorage（記録・設定）}};
  \node[side=impC, text width=3.6cm] (sb) at (2.1,-6.0)
    {\textgt{ネットの向こう}\\{\footnotesize Supabase（ログインだけ）}};
  \draw[arr] (user.south) -- node[lab,right]{見る・押す} (app.north);
  \draw[arr] (app.east) -- node[lab,above]{並べる} (comp.west);
  \draw[arr] (app.south) -- node[lab,right]{読む・書く} (ctx.north);
  \draw[arr] (ctx.east) -- node[lab,above]{計算を頼む} (utl.west);
  \draw[arr2] (as.north |- ctx.south) -- node[lab,left]{保存・読込} (as.north);
  \draw[arr2] (sb.north |- ctx.south) -- node[lab,right]{ログイン確認（\texttt{lib/}）} (sb.north);
  \draw[impA, dashed, line width=0.7pt] (typ.east |- app.west) -- (app.west);
  \draw[impA, dashed, line width=0.7pt] (typ.east |- ctx.west) -- (ctx.west);
\end{tikzpicture}
\end{center}
\figcap{図1　アプリ全体の地図（赤＝最重要、橙＝重要、灰＝参考）}


% =====================================================================
\section{定義}

この章では、このアプリを作っている技術を1つずつ定義し、「なぜそれを使っているのか」を書く。
最後に、コードを読むときに何度も出てくる React の用語をまとめる。

\begin{center}
\begin{tikzpicture}[font=\small,
  box/.style={draw=#1, fill=#1!7, line width=0.9pt, rounded corners=3pt,
    minimum height=0.95cm, align=center, inner sep=3pt},
  arr/.style={-{Stealth[length=2.2mm]}, line width=0.8pt, draw=black!60}]
  \node[box=impA, minimum width=14cm] (ts) at (0,0)
    {\textgt{TypeScript}：あなた（と AI）が書く言語　\texttt{.ts}／\texttt{.tsx}};
  \node[box=impA, minimum width=14cm] (re) at (0,-1.35)
    {\textgt{React}：画面を「部品」と「状態」で組み立てる考え方};
  \node[box=impB, minimum width=6.8cm] (rn) at (-3.6,-2.7)
    {\textgt{React Native}：部品をスマホの本物の部品にする};
  \node[box=impB, minimum width=6.8cm] (rw) at (3.6,-2.7)
    {\textgt{react-native-web}：部品をブラウザの部品にする};
  \node[box=impB, minimum width=14cm] (ex) at (0,-4.05)
    {\textgt{Expo（SDK 54）}：道具箱とビルドの肩代わり　＋　\textgt{expo-router}：ファイル＝画面};
  \node[box=impC, minimum width=4.4cm] (ios) at (-4.8,-5.5) {iPhone};
  \node[box=impC, minimum width=4.4cm] (and) at (0,-5.5) {Android};
  \node[box=impC, minimum width=4.4cm] (web) at (4.8,-5.5) {Web（Netlify で公開）};
  \draw[arr] (ts) -- (re);
  \draw[arr] (re.south -| rn) -- (rn);
  \draw[arr] (re.south -| rw) -- (rw);
  \draw[arr] (rn) -- (rn |- ex.north);
  \draw[arr] (rw) -- (rw |- ex.north);
  \draw[arr] (ex.south -| ios) -- (ios);
  \draw[arr] (ex.south -| and) -- (and);
  \draw[arr] (ex.south -| web) -- (web);
\end{tikzpicture}
\end{center}
\figcap{図2　技術の積み重ね（上で書いたものが、下へ行くほど「動くアプリ」に変わる）}

\subsection{言語：TypeScript}

\begin{kkbadgepoint}[定義]{TypeScript（タイプスクリプト）}
  JavaScript に「型（この変数には数字しか入らない、など）」を書けるようにした言語。
  書いた型が合わないと、\textgt{動かす前に}エラーで教えてくれる。
  最後は JavaScript に変換されて動く。このアプリのコードは、設定ファイルを除いてほぼすべて TypeScript。
\end{kkbadgepoint}

\textgt{なぜ使っているか。}
README に「\texttt{create-expo-app} で作った」とあり、その標準のひな形が TypeScript だから。\evCode
さらに \texttt{tsconfig.json} で \texttt{"strict": true}（型のチェックを最も厳しくする）が有効になっている。\evCode
型があると、AI が書いたコードの間違いも動かす前に見つかりやすい。\evGuess

\textgt{拡張子で言語を見分ける。}
\begin{center}
\begin{tblr}{
  colspec={Q[l,wd=2.6cm] X[l] Q[l,wd=5.2cm]},
  width=\linewidth,
  row{1}={bg=kkdark, fg=white, font=\small\gtfamily},
  hlines={black!15},
  cells={font=\small},
}
拡張子 & 中身 & このアプリの例 \\
\SetRow{bg=impAbg} \texttt{.tsx} & TypeScript ＋ JSX（画面の形を HTML のように書ける） & 画面と部品（\texttt{app/}・\texttt{components/}） \\
\SetRow{bg=impAbg} \texttt{.ts} & TypeScript（画面を持たない処理） & \texttt{types/}・\texttt{utils/}・\texttt{lib/} \\
\SetRow{bg=impBbg} \texttt{.json} & 設定データ（JSON 形式） & \texttt{package.json}・\texttt{app.json} \\
\SetRow{bg=impCbg} \texttt{.js} & JavaScript（型なし） & \texttt{eslint.config.js} \\
\SetRow{bg=impCbg} \texttt{.toml} & 設定データ（TOML 形式） & \texttt{netlify.toml} \\
\SetRow{bg=impCbg} \texttt{.md} & メモ（Markdown） & \texttt{log/}・\texttt{AGENTS.md} \\
\end{tblr}
\end{center}

TypeScript の書き方を、日付を作る小さな関数で見てみる。

\begin{codebox}{impB}{utils/date.ts}{7}
export function localDateKey(date: Date = new Date()): string {
  const y = date.getFullYear();
  const m = String(date.getMonth() + 1).padStart(2, '0');
  const d = String(date.getDate()).padStart(2, '0');
  return `${y}-${m}-${d}`;
}
\end{codebox}

\begin{pyaside}
  \texttt{def local\_date\_key(date: datetime = datetime.now()) -> str:} とほぼ同じ。
  \texttt{: Date} が引数の型、\texttt{): string} が戻り値の型。
  \texttt{const} は「あとで入れ直さない変数」。
  最後の行のバッククォート文字列は、Python の f 文字列 \texttt{f"\{y\}-\{m\}-\{d\}"} にあたる。
\end{pyaside}

\subsection{画面の仕組み：React と React Native}

\begin{kkbadgepoint}[定義]{React と React Native}
  \textgt{React} は、画面を「部品（コンポーネント）」の組み合わせで作り、
  部品が覚えている値（状態）が変わると、その部品を自動で描き直す仕組み。
  \textgt{React Native} は、React で書いた部品を、iPhone や Android の本物の部品に変えて表示する仕組み。
  \texttt{react-native-web} は、同じ部品をブラウザ用に変える。
\end{kkbadgepoint}

\textgt{なぜ使っているか。}
1つのコードで iPhone・Android・Web の3つに出せるから。
\texttt{app.json} に \texttt{ios}・\texttt{android}・\texttt{web} の3つの設定が並んでいることが、その使い方を示している。\evCode
ほかの方法（Swift と Kotlin で別々に作るなど）より、1人で作るときの手間が少ない。\evGuess

\subsection{開発の土台：Expo と expo-router}

\begin{kkbadgepoint}[定義]{Expo と expo-router}
  \textgt{Expo} は、React Native のアプリを作るための道具箱。
  保存・通知・画像などの部品（\texttt{expo-xxx}）と、アプリの組み立て（ビルド）を肩代わりする仕組みをまとめて提供する。
  このアプリは \textgt{SDK 54} を使う（\texttt{package.json} の \texttt{"expo": "\textasciitilde 54.0.34"}）。
  \textgt{expo-router} は、\texttt{app/} フォルダにファイルを置くと、そのファイルがそのまま1つの画面になる仕組み。
\end{kkbadgepoint}

expo-router には、ファイル名の決まりが2つある。
\begin{itemize}
  \item \texttt{\_layout.tsx}：同じフォルダの画面を包む「枠」。タブやヘッダーなど、共通の部分を書く。
  \item \texttt{(tabs)} のような括弧つきのフォルダ：画面をまとめるグループ。URL には出てこない。
\end{itemize}

\textgt{なぜ使っているか。}
README のとおり \texttt{create-expo-app} で作ったプロジェクトで、expo-router はその標準。\evCode
Xcode や Android Studio の細かい設定をしなくても、Expo Go というアプリで実機ですぐ試せる。\evGuess

アプリの入口は \texttt{package.json} に書いてある。

\begin{codebox}[language=plain]{impB}{package.json}{1}
{
  "name": "time_log_app",
  "main": "expo-router/entry",
  "version": "1.0.0",
  "scripts": {
    "start": "expo start",
    "android": "expo start --android",
    "ios": "expo start --ios",
    "web": "expo start --web",
    "lint": "expo lint"
  },
  "dependencies": {
    "@expo/vector-icons": "^15.0.3",
    "@react-native-async-storage/async-storage": "2.2.0",
    "@react-navigation/bottom-tabs": "^7.4.0",
    "@react-navigation/elements": "^2.6.3",
    "@react-navigation/native": "^7.1.8",
    "@supabase/supabase-js": "^2.110.0",
    "expo": "~54.0.34",
    "expo-blur": "~15.0.8",
\end{codebox}

3行目の \texttt{"main": "expo-router/entry"} が「起動したら expo-router に任せる」という意味。
5〜11行目の \texttt{scripts} は、\texttt{npm run web} などで呼べるコマンドの一覧。
12行目からの \texttt{dependencies} が、使っているライブラリの一覧（21行目以降も続く）。

\subsection{データの置き場所：AsyncStorage と Supabase}

\begin{kkbadgepoint}[定義]{AsyncStorage と Supabase}
  \textgt{AsyncStorage} は、端末の中にある「名前（キー）→ 文字列」の保存箱。
  このアプリは、学習記録・疲労・設定を JSON の文字列にして、ここに入れている。
  \textgt{Supabase} は、ネット上のサーバーをすぐ使えるサービス。
  このアプリでは\textgt{ログイン（メールとパスワード）だけ}に使っている。
\end{kkbadgepoint}

\begin{pyaside}
  開発ログ V1 によると、AsyncStorage は \texttt{json.dump}／\texttt{json.load} で端末のファイルに読み書きするのとほぼ同じ（ただし非同期）。\evLog
\end{pyaside}

\textgt{大事な事実。}コードの中で Supabase を呼んでいるのは \texttt{context/AuthContext.tsx} と \texttt{lib/supabase.ts} だけで、
使っているのはログイン用の \texttt{supabase.auth} だけ。学習記録をサーバーへ送るコードはない。\evCode
つまり、\textgt{別の端末で同じアカウントにログインしても、学習記録は引き継がれない。}

\textgt{なぜ使っているか。}
AsyncStorage はサーバーがいらず、端末の中だけで完結するので、最初に作るのに向いている。\evGuess
Supabase は、ログインの仕組み（パスワードの安全な保管など）を自分で作らずに済む。\evGuess

Supabase との接続は \texttt{lib/supabase.ts} にまとまっている。

\begin{codebox}{impB}{lib/supabase.ts}{12}
const supabaseUrl = process.env.EXPO_PUBLIC_SUPABASE_URL;
const supabaseAnonKey = process.env.EXPO_PUBLIC_SUPABASE_ANON_KEY;
\end{codebox}

接続先と鍵は、\texttt{mobile/.env.local} にある \texttt{EXPO\_PUBLIC\_SUPABASE\_URL} と
\texttt{EXPO\_PUBLIC\_SUPABASE\_ANON\_KEY} から読む。値そのものは秘密なので、このノートには載せていない。

\begin{codebox}{impB}{lib/supabase.ts}{34}
export const supabase = createClient(supabaseUrl, supabaseAnonKey, {
  auth: {
    storage: ssrSafeStorage,
    autoRefreshToken: true,
    persistSession: true,
    detectSessionInUrl: false,
  },
});
\end{codebox}

36行目では、ログイン状態の保存先にも AsyncStorage を使っている。
\texttt{ssrSafeStorage} という包み方をしているのは、Web 版を作るときにサーバー側で動く処理には \texttt{window} がなく、そのままだと落ちるから。\evCode
38行目の \texttt{persistSession: true} で、アプリを閉じてもログインしたままになる。

\subsection{公開：Netlify}

\begin{kkbadgepoint}[定義]{Netlify（ネットリファイ）}
  作った Web 版のファイル（HTML など）を置いて、インターネットで公開してくれるサービス。
  設定はリポジトリ直下の \texttt{netlify.toml} に書く。
\end{kkbadgepoint}

\begin{codebox}[language=plain]{impC}{netlify.toml}{1}
[build]
  # アプリ本体は mobile/ にある。command と publish は base からの相対パス
  base = "mobile"
  command = "npx expo export -p web"
  publish = "dist"

[build.environment]
  NODE_VERSION = "24"

[[redirects]]
  from = "/*"
  to = "/+not-found.html"
  status = 404
\end{codebox}

Netlify は \texttt{mobile/} に入り（3行目）、\texttt{npx expo export -p web} で Web 版を作って（4行目）、
できあがった \texttt{dist/} を公開する（5行目）。
\texttt{app.json} の \texttt{"output": "static"} により、Web 版は「サーバーのいらない、ただのファイルの集まり」として作られる。\evCode
そういうファイルを無料で手軽に公開できるのが Netlify を選ぶ理由として考えられる。\evGuess

\subsection{技術のまとめ}

\begin{center}
\begin{tblr}{
  colspec={Q[l,wd=3.1cm] X[l] X[l] Q[c,wd=1.7cm]},
  width=\linewidth,
  row{1}={bg=kkdark, fg=white, font=\small\gtfamily},
  hlines={black!15},
  cells={font=\small},
}
技術 & 役割 & なぜ & 根拠 \\
\SetRow{bg=impAbg} TypeScript & 書く言語 & ひな形の標準。型で早く間違いに気づく & コード／推測 \\
\SetRow{bg=impAbg} React & 部品と状態で画面を作る & React Native の土台 & 推測 \\
\SetRow{bg=impBbg} React Native & スマホの部品に変える & 1つのコードで3つに出せる & コード／推測 \\
\SetRow{bg=impBbg} Expo（SDK 54） & 道具箱とビルド & \texttt{create-expo-app} で作った & コード \\
\SetRow{bg=impBbg} expo-router & ファイル＝画面 & ひな形の標準 & コード \\
\SetRow{bg=impAbg} Context & 全画面でデータを共有 & 全画面で同じ値を使うため & ログ \\
\SetRow{bg=impAbg} AsyncStorage & 端末に保存 & サーバーなしで完結する & ログ／推測 \\
\SetRow{bg=impCbg} Supabase & ログインだけ & ログインを自作しなくて済む & 推測 \\
\SetRow{bg=impCbg} Netlify & Web 版の公開 & 静的なファイルを手軽に公開 & 推測 \\
\end{tblr}
\end{center}

画面の飾りのためのライブラリもある。キャラクターがふわふわ動くのは \texttt{react-native-reanimated}、
疲労が強いときのぼかしは \texttt{expo-blur}（開発ログ V3 に「Expo 公式で SDK 54 と互換が保証されている」とある \evLog）、
グラフは \texttt{react-native-svg}、CSV／JSON の書き出しは \texttt{expo-file-system} と \texttt{expo-sharing}。

\subsection{このあと何度も出てくる React の用語}

この先のコードを読むために、React の用語を先に定義しておく。

\begin{center}
\begin{tblr}{
  colspec={Q[l,wd=2.9cm] X[l] Q[l,wd=3.5cm] Q[l,wd=3.6cm]},
  width=\linewidth,
  row{1}={bg=kkdark, fg=white, font=\small\gtfamily},
  hlines={black!15},
  cells={font=\small},
}
用語 & 何か & Pythonでいうと & このアプリの例 \\
コンポーネント & 画面の部品を返す関数。名前は大文字で始める & 画面を組み立てて返す関数 & \texttt{TimerScreen()} \\
JSX & \texttt{<View>} や \texttt{<Text>} のように、画面の形を HTML 風に書く書き方 & 関数の中に直接書くテンプレート & \texttt{<Text>開始</Text>} \\
props & 部品に外から渡す値 & 関数の引数 & \texttt{visible}・\texttt{onDone} \\
state（\texttt{useState}） & 部品が覚えている値。専用の関数で変えると画面が描き直される & \texttt{self.seconds}。ただし代入ではなく \texttt{setSeconds(…)} で変える & \texttt{seconds}・\texttt{status} \\
\texttt{useEffect} & 描いた後に動かす処理。\texttt{[ ]} に書いた値が変わったときだけ動く & 値の変化を見張る関数 & 1秒ごとのタイマー \\
\texttt{useRef} & 描き直しても消えないが、変えても描き直さない入れ物 & ただのインスタンス変数 & \texttt{currentSessionIdRef} \\
Context と Provider & アプリ全体で共有するデータ。Provider で囲んだ内側の画面だけが使える & Flask の \texttt{g}（開発ログ V1） & \texttt{TimeProvider} \\
フック & \texttt{use} で始まる関数の総称。自分で作ることもできる & なし & \texttt{useTime()} \\
\end{tblr}
\end{center}

Context は「どこで囲んでいるか」が大事なので、図3に描いた。外側の箱ほど広い範囲から使える。

\begin{center}
\begin{tikzpicture}[font=\footnotesize,
  lab/.style={anchor=north west, inner sep=1pt, align=left}]
  \draw[impB, fill=impB!4, line width=0.9pt, rounded corners=3pt] (0,0) rectangle (16.4,-7.0);
  \node[lab] at (0.2,-0.15) {\textcolor{impB}{\textgt{\texttt{AuthProvider}}}：ログインしているか\quad（\texttt{app/\_layout.tsx}）};
  \draw[impC, fill=impCbg, line width=0.9pt, rounded corners=3pt] (0.35,-0.75) rectangle (16.05,-6.65);
  \node[lab] at (0.55,-0.9) {\textcolor{impC}{\textgt{\texttt{SettingsProvider}}}：テーマ・換算比率};
  \draw[black!40, fill=white, line width=0.7pt, rounded corners=3pt] (0.75,-1.55) rectangle (3.65,-6.25);
  \node[align=center] at (2.2,-3.9) {\textgt{ログイン画面}\\\texttt{login.tsx}\\[4pt]{\scriptsize 学習記録の箱の}\\{\scriptsize 外側にある}};
  \draw[black!40, fill=white, line width=0.7pt, rounded corners=3pt] (4.0,-1.55) rectangle (15.7,-6.25);
  \node[lab] at (4.2,-1.7) {\textgt{タブ画面のグループ}\quad（\texttt{app/(tabs)/\_layout.tsx}）};
  \draw[impA, fill=impA!4, line width=0.9pt, rounded corners=3pt] (4.3,-2.3) rectangle (15.4,-5.95);
  \node[lab] at (4.5,-2.45) {\textcolor{impA}{\textgt{\texttt{TimeProvider}}}：学習時間（累計と日ごと）};
  \draw[impB, fill=impB!5, line width=0.9pt, rounded corners=3pt] (4.6,-3.0) rectangle (15.1,-5.65);
  \node[lab] at (4.8,-3.15) {\textcolor{impB}{\textgt{\texttt{FatigueProvider}}}：疲労（予定とスコア）};
  \draw[impA, fill=impA!8, line width=0.9pt, rounded corners=3pt] (4.9,-3.7) rectangle (14.8,-5.35);
  \node[lab] at (5.1,-3.85) {\textcolor{impA}{\textgt{\texttt{StudyLogProvider}}}：前後チェックの記録};
  \foreach \x/\t in {6.05/ホーム, 7.95/タイマー, 9.85/カレンダー, 11.75/疲労, 13.65/設定} {
    \node[draw=black!50, fill=white, rounded corners=2pt, minimum width=1.7cm, minimum height=0.6cm] at (\x,-4.8) {\t};
  }
\end{tikzpicture}
\end{center}
\figcap{図3　Provider（共有データの箱）の入れ子。タブの5画面は、すべての箱の内側にある}

図3の入れ子は、コードではタグの入れ子としてそのまま書かれている。

\begin{codebox}{impA}{app/(tabs)/\_layout.tsx}{12}
export default function TabLayout() {
  const colorScheme = useColorScheme();

  return (
    <TimeProvider>
      <FatigueProvider>
        <StudyLogProvider>
          <Tabs
            screenOptions={{
              tabBarActiveTintColor: Colors[colorScheme ?? 'light'].tint,
              headerShown: false,
              tabBarButton: HapticTab,
            }}
\end{codebox}

16〜18行目で3つの Provider が順に包み、その内側の \texttt{<Tabs>} にタブの5画面がぶら下がる。
Context を「学習時間」「疲労」「前後チェック」に分けているのは、1つにまとめると直すときに互いに影響するから、と開発ログ V3 に書かれている。\evLog


% =====================================================================
\section{要点}

この章では、フォルダの地図 → データの形 → 1つの機能の流れ → 細部、の順に読んでいく。

\subsection{フォルダの地図}

\begin{kkbadgepoint}[要点]{最初に読むのは赤の場所だけでいい}
  \texttt{mobile/} には約50個のファイルがあるが、このアプリ独自の中心は
  \texttt{app/(tabs)/}・\texttt{context/}・\texttt{types/} の3か所。
  \texttt{components/} の一部や \texttt{hooks/} の一部は、Expo が最初に用意した雛形がそのまま残っているだけで、
  4つは今どこからも使われていない。自動生成の \texttt{node\_modules/} や \texttt{dist/} は読む必要がない。
\end{kkbadgepoint}

\textgt{\texttt{mobile/}（アプリ本体）の中身}

\begin{longtblr}[entry=none,label=none]{
  colspec={Q[l,wd=4.3cm] X[l] Q[c,wd=1.8cm] Q[c,wd=1.6cm]},
  width=\linewidth,
  rowhead=1,
  row{1}={bg=kkdark, fg=white, font=\small\gtfamily},
  hlines={black!15},
  cells={font=\small},
}
場所 & 役割（一言） & 種類 & 重要度 \\
\SetRow{bg=impAbg} \texttt{app/(tabs)/\_layout.tsx} & 下のタブ5つと、共有データの箱（Provider）の配置 & 独自 & \lvA \\
\SetRow{bg=impAbg} \texttt{app/(tabs)/index.tsx} & ホーム（キャラ・累計・今日の時間） & 独自 & \lvA \\
\SetRow{bg=impAbg} \texttt{app/(tabs)/timer.tsx} & タイマー（開始・停止と前後チェック） & 独自 & \lvA \\
\SetRow{bg=impBbg} \texttt{app/(tabs)/history.tsx} & カレンダー（日ごとの学習時間） & 独自 & \lvB \\
\SetRow{bg=impBbg} \texttt{app/(tabs)/fatigue.tsx} & 疲労チェック（予定の入力とバッテリー） & 独自 & \lvB \\
\SetRow{bg=impBbg} \texttt{app/(tabs)/settings.tsx} & 設定（テーマ・換算・グラフ・書き出し・削除） & 独自 & \lvB \\
\SetRow{bg=impBbg} \texttt{app/\_layout.tsx} & 一番外側の枠。ログインしていなければログイン画面へ送る & 独自 & \lvB \\
\SetRow{bg=impBbg} \texttt{app/login.tsx} & ログイン・新規登録・ゲストで使う & 独自 & \lvB \\
\SetRow{bg=impCbg} \texttt{app/modal.tsx} & 雛形の見本画面。どこからも開かれていない & 雛形（未使用） & \lvC \\
\SetRow{bg=impAbg} \texttt{context/} & アプリ全体で共有するデータと、その保存・読込（5つ） & 独自 & \lvA \\
\SetRow{bg=impAbg} \texttt{types/} & データの形（型）の定義だけ（2つ） & 独自 & \lvA \\
\SetRow{bg=impBbg} \texttt{utils/} & 計算・変換の関数。画面を持たない（7つ） & 独自 & \lvB \\
\SetRow{bg=impBbg} \texttt{hooks/use-fatigue.ts} & 疲労の状態を計算して画面に渡す & 独自 & \lvB \\
\SetRow{bg=impCbg} \texttt{hooks/} の残り3つ & ダーク／ライトの判定と色選び & 雛形 & \lvC \\
\SetRow{bg=impBbg} \texttt{components/} の独自8つ & 疲労バッテリー・ぼかし・反応時間テスト・点数選び・グラフ2つ・前後チェック2つ & 独自 & \lvB \\
\SetRow{bg=impCbg} \texttt{components/} の雛形4つ（使用中） & \texttt{themed-text}・\texttt{themed-view}・\texttt{haptic-tab}・\texttt{ui/icon-symbol} & 雛形 & \lvC \\
\SetRow{bg=impDbg} \texttt{components/} の雛形4つ（未使用） & \texttt{hello-wave}・\texttt{parallax-scroll-view}・\texttt{external-link}・\texttt{ui/collapsible} & 雛形（未使用） & \lvD \\
\SetRow{bg=impBbg} \texttt{lib/supabase.ts} & Supabase（ログイン用サーバー）への接続 & 独自 & \lvB \\
\SetRow{bg=impCbg} \texttt{constants/theme.ts} & 色とフォント。\texttt{Colors} は雛形、\texttt{AppColors}（緑）は独自 & 混在 & \lvC \\
\SetRow{bg=impBbg} \texttt{package.json} & 使うライブラリとコマンドの一覧 & 設定 & \lvB \\
\SetRow{bg=impCbg} \texttt{app.json} & アプリ名・アイコン・対応 OS など Expo の設定 & 設定 & \lvC \\
\SetRow{bg=impCbg} \texttt{tsconfig.json} & TypeScript の設定（\texttt{strict}、\texttt{@/} の別名） & 設定 & \lvC \\
\SetRow{bg=impCbg} \texttt{eslint.config.js} & 書き方のチェックの設定 & 雛形 & \lvC \\
\SetRow{bg=impCbg} \texttt{assets/} & アイコンや起動画面の画像 & 素材 & \lvC \\
\SetRow{bg=impCbg} \texttt{README.md} & Expo の雛形の説明書のまま & 雛形 & \lvC \\
\SetRow{bg=impDbg} \texttt{node\_modules/}・\texttt{package-lock.json} & インストールしたライブラリの本体と、その版の記録 & 自動生成 & \lvD \\
\SetRow{bg=impDbg} \texttt{.expo/}・\texttt{dist/}・\texttt{expo-env.d.ts} & 開発中の一時ファイル・Web 版の完成品・型の補助 & 自動生成 & \lvD \\
\SetRow{bg=impDbg} \texttt{.env.local} & Supabase の接続先などの秘密の値。開かない & 秘密 & \lvD \\
\end{longtblr}

\texttt{@/} は \texttt{tsconfig.json} で決めた「\texttt{mobile/} を指す別名」。
\texttt{import \{ useTime \} from '@/context/TimeContext'} は「\texttt{mobile/context/TimeContext.tsx} から読み込む」という意味になる。

\textgt{リポジトリ直下（アプリの外）}

\begin{center}
\begin{tblr}{
  colspec={Q[l,wd=4.3cm] X[l] Q[c,wd=1.6cm]},
  width=\linewidth,
  row{1}={bg=kkdark, fg=white, font=\small\gtfamily},
  hlines={black!15},
  cells={font=\small},
}
場所 & 役割（一言） & 重要度 \\
\SetRow{bg=impBbg} \texttt{log/} & 開発ログ V1〜V3。「なぜ」の手がかりがある & \lvB \\
\SetRow{bg=impCbg} \texttt{CLAUDE.md}・\texttt{AGENTS.md} & AI への指示（フォルダ構成と安全ルール） & \lvC \\
\SetRow{bg=impCbg} \texttt{netlify.toml} & Web 版を公開する設定 & \lvC \\
\SetRow{bg=impCbg} \texttt{.gitignore} & Git に入れないファイルの一覧 & \lvC \\
\end{tblr}
\end{center}

\subsection{データの形（types/）}

\begin{kkbadgepoint}[要点]{コードを読む前に「何を覚えているアプリか」を知る}
  \texttt{types/} には処理が1行もなく、「データがどんな項目でできているか」だけが書いてある。
  ここを先に見ると、画面や context のコードが「この項目を作っている」「この項目を足している」と読めるようになる。
  学習時間そのもの（累計と日ごと）は \texttt{types/} ではなく、\texttt{context/TimeContext.tsx} の中で形を決めている。
\end{kkbadgepoint}

\textgt{TypeScript の型の読み方}（このあとのコードに出てくるもの）
\begin{itemize}
  \item \texttt{interface 名前 \{ … \}}：項目の一覧。Python の \texttt{TypedDict} や \texttt{dataclass} に近い。
  \item \texttt{subject?: string}：\texttt{?} は「あってもなくてもよい」。Python の \texttt{Optional[str]}。
  \item \texttt{'pre' | 'post'}：\texttt{|} は「どれか1つ」。この場合は2つの文字列のどちらか。Python の \texttt{Literal['pre', 'post']}。
  \item \texttt{number | null}：数字か、何もない（\texttt{null}）か。
  \item \texttt{Record<string, number>}：キーが文字列、値が数字の辞書。Python の \texttt{dict[str, int]}。
  \item \texttt{ReactionTrial[]}：\texttt{[]} は配列。Python の \texttt{list[ReactionTrial]}。
\end{itemize}

\begin{center}
\begin{tikzpicture}[font=\footnotesize,
  ent/.style={draw=#1, fill=white, line width=0.9pt, rounded corners=2pt, inner sep=4pt},
  arr/.style={-{Stealth[length=2.2mm]}, line width=0.8pt, draw=black!60},
  lab/.style={font=\scriptsize, fill=white, inner sep=1.2pt}]
  \node[ent=impA] (ss) at (0,0) {\begin{tabular}{@{}l@{}}
    \textcolor{impA}{\textgt{StudySession}}（1回の学習）\\\hline
    \texttt{id}：この記録の番号\\
    \texttt{subject?}：科目\\
    \texttt{startedAt}／\texttt{endedAt?}：開始・終了の日時\\
    \texttt{durationMinutes}：学習した分数\\
    \texttt{reward}：面白さ・重要度・締切\\
    \texttt{subjectivePreId?} ほか：チェックの番号
  \end{tabular}};
  \node[ent=impB] (se) at (7.9,1.3) {\begin{tabular}{@{}l@{}}
    \textcolor{impB}{\textgt{SubjectiveEntry}}（自分で感じた疲れ）\\\hline
    \texttt{sessionId?}：どの学習のものか\\
    \texttt{timing}：\texttt{'pre'}（前）か \texttt{'post'}（後）\\
    \texttt{score}：疲労感 1〜5\\
    \texttt{note?}：メモ
  \end{tabular}};
  \node[ent=impB] (rr) at (7.9,-2.1) {\begin{tabular}{@{}l@{}}
    \textcolor{impB}{\textgt{ReactionTimeResult}}（反応時間テスト）\\\hline
    \texttt{sessionId?}：どの学習のものか\\
    \texttt{timing}：前か後\\
    \texttt{trials}：1回ずつの結果\\
    \texttt{avgReactionMs}：平均の反応時間\\
    \texttt{accuracy}：正確さ 0〜1
  \end{tabular}};
  \node[ent=impC] (ra) at (0,-4.3) {\begin{tabular}{@{}l@{}}
    \textcolor{impC}{\textgt{RewardAppraisal}}（やる気の材料）\\\hline
    \texttt{interest}：面白さ 1〜5\\
    \texttt{importance}：重要度 1〜5\\
    \texttt{hasDeadline}／\texttt{deadlineUrgency?}：締切
  \end{tabular}};
  \node[ent=impC] (rt) at (7.9,-5.4) {\begin{tabular}{@{}l@{}}
    \textcolor{impC}{\textgt{ReactionTrial}}（1回分）\\\hline
    \texttt{delayMs}：光るまでの待ち時間\\
    \texttt{reactionMs}：反応時間（\texttt{null} は失敗）
  \end{tabular}};
  \draw[arr] (se.west) -- node[lab,above,sloped]{\texttt{sessionId} で結ぶ} (ss.east);
  \draw[arr] (rr.west) -- node[lab,below,sloped]{\texttt{sessionId} で結ぶ} (ss.east);
  \draw[arr] (ss.south) -- node[lab,right]{\texttt{reward} として中に持つ} (ra.north);
  \draw[arr] (rr.south) -- node[lab,right]{\texttt{trials} に5回分} (rt.north);
\end{tikzpicture}
\end{center}
\figcap{図4　前後チェックの記録の形（\texttt{types/studyLog.ts}）。矢印は「どれとつながるか」}

\begin{codebox}{impA}{types/studyLog.ts}{1}
export const LOCAL_USER_ID = 'local-user-001';

export interface RewardAppraisal {
  interest: number; // 1-5 面白さ
  importance: number; // 1-5 重要度
  hasDeadline: boolean;
  deadlineUrgency?: number; // 1-5、hasDeadline=trueの時のみ
}

export interface StudySession {
  id: string;
  userId: string;
  subject?: string;
  startedAt: string; // ISO
  endedAt?: string; // ISO
  durationMinutes: number;
  reward: RewardAppraisal | null;
  subjectivePreId?: string;
  subjectivePostId?: string;
  reactionPreId?: string;
  reactionPostId?: string;
  /** 将来の「予定」欄用。現時点ではUIなし、欠損許容。 */
  plannedFor?: string; // YYYY-MM-DD
}

export type SubjectiveTiming = 'pre' | 'post';

export interface SubjectiveEntry {
  id: string;
  userId: string;
  sessionId?: string;
  timing: SubjectiveTiming;
  score: number; // 1-5（疲労感・高いほど疲れている）
  note?: string;
  timestamp: string; // ISO
}

export interface ReactionTrial {
  delayMs: number;
  reactionMs: number | null; // null = 反応なし/フライング
}
\end{codebox}

1行目の \texttt{LOCAL\_USER\_ID} は、全部の記録に入れている仮のユーザー番号。
ログインしても、記録はこの同じ番号で端末に保存される。
14行目の \texttt{// ISO} は、日時を \texttt{"2026-11-17T05:30:00.000Z"} のような文字列で持つという意味（末尾の \texttt{Z} は世界標準時）。

疲労の予定は、もう1つのファイルで形を決めている。

\begin{codebox}{impB}{types/fatigue.ts}{1}
export type FatigueAxis = 'cognitive' | 'physical';
export type FatigueWeight = 1 | 2 | 5;
export type FatigueLevel = 'ok' | 'caution' | 'danger';

export interface ScheduleTag {
  id: string;
  name: string;
  axis: FatigueAxis;
  weight: FatigueWeight;
  emoji: string;
}

/**
 * 正規化された予定エントリ。
 * v3: source='manual'（手動入力）のみ。
 * v4: source='calendar'（Googleカレンダー等）を追加予定。
 * calcFatigueScore() はどちらのソースも同一ロジックで処理できる。
 */
export interface ScheduleEntry {
  id: string;
  tagId: string;
  date: string; // YYYY-MM-DD
  durationMinutes: number;
  source: 'manual' | 'calendar';
  title?: string; // v4でカレンダーイベント名を上書き
  externalId?: string; // v4でカレンダーイベントID
}
\end{codebox}

\texttt{ScheduleTag} は「部活」「授業」のような予定の種類で、\texttt{axis} が身体（\texttt{physical}）か頭（\texttt{cognitive}）か、
\texttt{weight} がその重さ（1・2・5のどれか）。\texttt{ScheduleEntry} は「その種類の予定を、何日に何分やったか」。
21行目の \texttt{tagId} で \texttt{ScheduleTag} とつながる。

最後に、学習時間そのものの形。これは \texttt{TimeContext.tsx} の中にある。

\begin{codebox}{impA}{context/TimeContext.tsx}{6}
const STORAGE_KEY = 'timeLogData';

type SavedData = {
  totalMinutes: number;
  history: Record<string, number>;
};
\end{codebox}

\texttt{totalMinutes} は累計の分数、\texttt{history} は「日付 → その日の分数」の辞書。
たとえば \texttt{\{ "2026-11-16": 50, "2026-11-17": 45 \}} のようになる。

\textgt{端末の保存箱（AsyncStorage）の一覧}

\begin{center}
\begin{tblr}{
  colspec={Q[l,wd=3.6cm] X[l] Q[l,wd=3.6cm] Q[c,wd=1.6cm]},
  width=\linewidth,
  row{1}={bg=kkdark, fg=white, font=\small\gtfamily},
  hlines={black!15},
  cells={font=\small},
}
キー（箱の名前） & 中身 & 担当 & 重要度 \\
\SetRow{bg=impAbg} \texttt{timeLogData} & 累計の分数と、日ごとの分数 & \texttt{TimeContext.tsx} & \lvA \\
\SetRow{bg=impAbg} \texttt{studyLogData\_v1} & 学習1回ごとの記録・前後の疲労感・反応時間 & \texttt{StudyLogContext.tsx} & \lvA \\
\SetRow{bg=impBbg} \texttt{fatigueData\_v3} & 予定の種類と、日ごとの予定 & \texttt{FatigueContext.tsx} & \lvB \\
\SetRow{bg=impCbg} \texttt{appSettings} & テーマと換算比率 & \texttt{SettingsContext.tsx} & \lvC \\
\SetRow{bg=impCbg} \texttt{authGuestMode} & ゲストで使っているかどうか & \texttt{AuthContext.tsx} & \lvC \\
\end{tblr}
\end{center}

\subsection{1つの機能を縦に追う：タイマー → 保存 → ホーム表示}

ここがこのノートの中心。「タイマーで開始して、停止すると、ホームの『今日の学習時間』が増える」までを、
図5の番号①〜⑬の順に、元のコードで追う。

\begin{center}
\begin{tikzpicture}[font=\scriptsize,
  head/.style={draw=#1, fill=#1!10, line width=0.8pt, rounded corners=2pt,
    minimum width=2.6cm, minimum height=1.35cm, align=center, anchor=north, inner sep=2pt},
  msg/.style={-{Stealth[length=2mm]}, line width=0.7pt, draw=black!75},
  ret/.style={-{Stealth[length=2mm]}, line width=0.7pt, draw=black!60, dashed},
  ml/.style={fill=white, inner sep=1.2pt, align=center, above}]
  \node[head=impC] (hu) at (0,0) {\footnotesize\textgt{あなた}};
  \node[head=impA] (ht) at (2.95,0) {\footnotesize\textgt{タイマー画面}\\\texttt{timer.tsx}};
  \node[head=impB] (hm) at (5.9,0) {\footnotesize\textgt{前後チェック}\\\texttt{PreSessionModal}\\\texttt{PostSessionModal}};
  \node[head=impA] (hc) at (8.85,0) {\footnotesize\textgt{共有データ}\\\texttt{StudyLogContext}\\\texttt{TimeContext}};
  \node[head=impB] (hs) at (11.8,0) {\footnotesize\textgt{端末の保存箱}\\AsyncStorage};
  \node[head=impA] (hh) at (14.75,0) {\footnotesize\textgt{ホーム画面}\\\texttt{index.tsx}};
  \foreach \n in {hu,ht,hm,hc,hs,hh} { \draw[black!30, dashed] (\n.south) -- (\n.south |- 0,-12.6); }
  % 開始
  \draw[msg] (0,-1.9) -- node[ml]{①「開始」を押す} (2.95,-1.9);
  \draw[msg] (2.95,-2.75) -- node[ml]{② 前チェックを開く} (5.9,-2.75);
  \draw[msg] (5.9,-3.6) -- node[ml]{③ 入力を渡す（\texttt{onDone}）} (2.95,-3.6);
  \draw[msg] (2.95,-4.45) -- node[ml]{④ \texttt{startSession()} で記録を作る} (8.85,-4.45);
  \draw[msg] (8.85,-5.3) -- node[ml]{⑤ \texttt{studyLogData\_v1} に保存} (11.8,-5.3);
  \node[draw=impA, fill=impAbg, rounded corners=2pt, align=center, inner sep=2pt] at (2.95,-6.15)
    {⑥ 1秒ごとに \texttt{seconds} を +1};
  % 停止
  \draw[msg] (0,-7.0) -- node[ml]{⑦「停止」を押す} (2.95,-7.0);
  \draw[msg] (2.95,-7.85) -- node[ml]{⑧ 後チェックを開く} (5.9,-7.85);
  \draw[msg] (5.9,-8.7) -- node[ml]{⑨ 入力を渡す（\texttt{onDone}）} (2.95,-8.7);
  \draw[msg] (2.95,-9.55) -- node[ml]{⑩ \texttt{addMinutes(分)} と \texttt{finishSession()}} (8.85,-9.55);
  \draw[msg] (8.85,-10.4) -- node[ml]{⑪ \texttt{timeLogData} に保存} (11.8,-10.4);
  \draw[msg] (8.85,-11.25) -- node[ml]{⑫ 値が変わったので描き直す（\texttt{useTime()}）} (14.75,-11.25);
  \draw[ret] (14.75,-12.1) -- node[ml]{⑬「今日の学習時間」が増えて見える} (0,-12.1);
\end{tikzpicture}
\end{center}
\figcap{図5　タイマーを使ってから、ホームの数字が増えるまで（実線＝呼び出し、点線＝あなたに見える結果）}

\subsubsection{①② 「開始」を押すと、前チェックが開く}

\begin{codebox}{impA}{app/(tabs)/timer.tsx}{132}
        <View style={styles.buttonArea}>
          {status === 'idle' && (
            <TouchableOpacity
              style={[styles.btn, styles.btnPrimary]}
              onPress={() => setPreModalVisible(true)}
              activeOpacity={0.8}
            >
              <Text style={styles.btnTextLight}>開始</Text>
            </TouchableOpacity>
          )}
\end{codebox}

133行目の \texttt{status === 'idle' \&\&} は「状態が \texttt{idle}（止まっている）のときだけ、このボタンを表示する」という書き方。
136行目の \texttt{onPress} が押したときの処理で、\texttt{setPreModalVisible(true)} で前チェックの画面を「見える」にする。
\textgt{このボタンはまだタイマーを動かさない。}動かすのは③のあと。

タイマー画面が覚えている値（state）は、関数の最初にまとめて書いてある。

\begin{codebox}{impA}{app/(tabs)/timer.tsx}{23}
export default function TimerScreen() {
  const colorScheme = useColorScheme();
  const isDark = colorScheme === 'dark';

  const { addMinutes } = useTime();
  const { startSession, finishSession, addSubjectiveEntry, addReactionResult } = useStudyLogContext();
  const [status, setStatus] = useState<TimerStatus>('idle');
  const [seconds, setSeconds] = useState(0);
  const intervalRef = useRef<ReturnType<typeof setInterval> | null>(null);
  const [preModalVisible, setPreModalVisible] = useState(false);
  const [postModalVisible, setPostModalVisible] = useState(false);
  const currentSessionIdRef = useRef<string | null>(null);
\end{codebox}

27〜28行目で、共有データ（Context）から使う関数を取り出している。
29〜33行目が state で、\texttt{status} は \texttt{'idle'}（停止）・\texttt{'running'}（計測中）・\texttt{'paused'}（一時停止）のどれか。
34行目の \texttt{currentSessionIdRef} は「今計っている学習の番号」で、画面の見た目に関係しないので state ではなく \texttt{useRef} にしている。

\subsubsection{③④⑤ 前チェックの入力を受け取り、記録を作って保存する}

\begin{codebox}{impA}{app/(tabs)/timer.tsx}{61}
  const handlePreSessionDone: ComponentProps<typeof PreSessionModal>['onDone'] = input => {
    const sessionId = startSession({ subject: input.subject, reward: input.reward });
    if (input.subjectiveScore != null) {
      addSubjectiveEntry({
        sessionId,
        timing: 'pre',
        score: input.subjectiveScore,
        note: input.subjectiveNote,
      });
    }
    if (input.reaction) {
      addReactionResult({
        sessionId,
        timing: 'pre',
        trials: input.reaction.trials,
        avgReactionMs: input.reaction.avgReactionMs,
        accuracy: input.reaction.accuracy,
      });
    }
    currentSessionIdRef.current = sessionId;
    setPreModalVisible(false);
    setStatus('running');
  };
\end{codebox}

前チェックで「完了」を押すと、この関数が呼ばれる（③）。
62行目で \texttt{startSession} を呼び、新しい学習の記録を作ってその番号をもらう（④）。
63〜79行目は、疲労感や反応時間を入力した場合だけ、それも記録する（入力は任意なので \texttt{if} で確かめている）。
82行目の \texttt{setStatus('running')} で、ここで初めてタイマーが動き出す。

\begin{codebox}{impA}{context/StudyLogContext.tsx}{107}
  const startSession = useCallback(
    (input: { subject?: string; reward: RewardAppraisal | null; plannedFor?: string }) => {
      const id = genId();
      const session: StudySession = {
        id,
        userId: LOCAL_USER_ID,
        subject: input.subject,
        startedAt: new Date().toISOString(),
        durationMinutes: 0,
        reward: input.reward,
        plannedFor: input.plannedFor,
      };
      setSessions(prev => {
        const next = [...prev, session];
        persist({ sessions: next });
        return next;
      });
      return id;
    },
    [persist]
  );
\end{codebox}

110〜118行目で、図4の \texttt{StudySession} の形どおりに記録を作る。この時点では \texttt{durationMinutes} は 0。
120行目で、今までの記録の配列の後ろに新しい記録を足した配列を作り、
121行目の \texttt{persist} で AsyncStorage の \texttt{studyLogData\_v1} に保存する（⑤）。

\begin{pyaside}
  120行目の \texttt{[.{}.{}.prev, session]} は \texttt{[*prev, session]}、
  つまり「\texttt{prev} の中身を全部並べて、最後に \texttt{session} を足した新しいリスト」。
  \texttt{prev.append(session)} のように元のリストを書き換えないのは、
  React が「新しい値になった」と気づいて描き直すため。
\end{pyaside}

\subsubsection{⑥ 1秒ごとに数える}

\begin{codebox}{impA}{app/(tabs)/timer.tsx}{36}
  useEffect(() => {
    if (status === 'running') {
      intervalRef.current = setInterval(() => {
        setSeconds(prev => prev + 1);
      }, 1000);
    } else {
      if (intervalRef.current) {
        clearInterval(intervalRef.current);
        intervalRef.current = null;
      }
    }

    return () => {
      if (intervalRef.current) {
        clearInterval(intervalRef.current);
        intervalRef.current = null;
      }
    };
  }, [status]);
\end{codebox}

54行目の \texttt{[status]} により、この処理は \texttt{status} が変わったときだけ動く。
\texttt{'running'} になったら、38〜40行目の \texttt{setInterval} で 1000 ミリ秒（1秒）ごとに \texttt{seconds} を 1 増やす。
一時停止や停止で \texttt{'running'} でなくなったら、43行目の \texttt{clearInterval} で止める。
48〜53行目の \texttt{return} の中身は「後片付け」で、画面を離れたときにもタイマーを止める。

\begin{pyaside}
  開発ログ V1 によると、\texttt{setInterval}／\texttt{clearInterval} は \texttt{threading.Timer} にあたる。\evLog
  39行目の \texttt{prev => prev + 1} は \texttt{lambda prev: prev + 1}。
\end{pyaside}

\subsubsection{⑦⑧⑨⑩ 「停止」を押し、後チェックのあとで分数を足す}

\begin{codebox}{impA}{app/(tabs)/timer.tsx}{56}
  const handleStop = () => {
    setStatus('idle');
    setPostModalVisible(true);
  };
\end{codebox}

停止ボタンは、タイマーを止めて（57行目）後チェックを開く（58行目）だけ。まだ何も保存しない。

\begin{codebox}{impA}{app/(tabs)/timer.tsx}{85}
  const handlePostSessionDone: ComponentProps<typeof PostSessionModal>['onDone'] = input => {
    const elapsedMinutes = Math.round(seconds / 60);
    addMinutes(elapsedMinutes);

    const sessionId = currentSessionIdRef.current;
    if (sessionId) {
      finishSession(sessionId, elapsedMinutes);
\end{codebox}

後チェックで「完了」を押すと、この関数が呼ばれる（⑨）。
86行目で秒を分に直す。\texttt{Math.round} は四捨五入なので、29秒までは 0 分、30秒からは 1 分になる。
87行目の \texttt{addMinutes} で \texttt{TimeContext} に分数を足し、
91行目の \texttt{finishSession} で \texttt{StudySession} に終了時刻と分数を書き込む（⑩）。
このあと92〜108行目は、③のときと同じように \texttt{'post'}（後）の疲労感と反応時間を記録する。

\begin{codebox}{impA}{app/(tabs)/timer.tsx}{109}
    }
    currentSessionIdRef.current = null;
    setPostModalVisible(false);
    setSeconds(0);
  };
\end{codebox}

最後に、学習の番号を空にして、後チェックを閉じ、秒を 0 に戻す。

\subsubsection{⑪ TimeContext が分数を足して保存する}

\begin{codebox}{impA}{context/TimeContext.tsx}{59}
  useEffect(() => {
    if (!loaded) return;
    const data: SavedData = { totalMinutes, history };
    AsyncStorage.setItem(STORAGE_KEY, JSON.stringify(data)).catch(() => {});
  }, [totalMinutes, history, loaded]);

  const addMinutes = (minutes: number) => {
    if (minutes <= 0) return;
    const today = todayKey();
    setTotalMinutes(prev => prev + minutes);
    setHistory(prev => ({
      ...prev,
      [today]: (prev[today] ?? 0) + minutes,
    }));
  };
\end{codebox}

\textgt{読む順番は下から。}⑩で呼ばれるのは 65行目の \texttt{addMinutes}。
\begin{itemize}
  \item 66行目：0 分以下なら何もしない。
  \item 67行目：今日の日付 \texttt{"2026-11-17"} のような文字列を作る（\texttt{utils/date.ts} の関数）。
  \item 68行目：累計に足す。
  \item 69〜72行目：日ごとの辞書の、今日の分に足す。今日の分がまだなければ 0 から足す（\texttt{?? 0}）。
\end{itemize}
68〜72行目で state が変わると、59〜63行目の \texttt{useEffect} が「\texttt{totalMinutes} と \texttt{history} が変わった」と気づいて動き、
JSON の文字列にして \texttt{timeLogData} に保存する（⑪）。
\textgt{「足す」処理と「保存する」処理が分かれていて、保存は state の変化をきっかけに自動で起きる}のがポイント。

\begin{pyaside}
  62行目の \texttt{JSON.stringify(data)} は \texttt{json.dumps(data)}。
  70〜71行目の \texttt{\{ .{}.{}.prev, [today]: … \}} は \texttt{\{**prev, today: …\}}。
  71行目の \texttt{prev[today] ?? 0} は \texttt{prev.get(today, 0)} とほぼ同じ。
\end{pyaside}

\subsubsection{⑫⑬ ホーム画面が描き直される}

\begin{codebox}{impA}{context/TimeContext.tsx}{29}
export function TimeProvider({ children }: { children: ReactNode }) {
  const [totalMinutes, setTotalMinutes] = useState(0);
  const [history, setHistory] = useState<Record<string, number>>({});
  const [loaded, setLoaded] = useState(false);

  const todayMinutes = history[todayKey()] ?? 0;
\end{codebox}

\texttt{todayMinutes}（今日の分数）は保存されている値ではなく、34行目で \texttt{history} から毎回計算している。
\texttt{history} が変わると \texttt{TimeProvider} が描き直され、この値も新しくなる。
そして、\texttt{useTime()} でこの値を使っているすべての画面に「値が変わった」と伝わる（⑫）。

\begin{codebox}{impA}{app/(tabs)/index.tsx}{61}
export default function HomeScreen() {
  const colorScheme = useColorScheme();
  const isDark = colorScheme === 'dark';
  const { totalMinutes, todayMinutes } = useTime();
  const { fatigueState } = useFatigue();
\end{codebox}

\begin{codebox}{impA}{app/(tabs)/index.tsx}{170}
          <View style={styles.statItem}>
            <ThemedText style={styles.statLabel}>今日の学習時間</ThemedText>
            <View style={styles.statValueRow}>
              <ThemedText style={styles.statValueBig}>{todayMinutes}</ThemedText>
              <ThemedText style={styles.statUnit}>分</ThemedText>
            </View>
          </View>
\end{codebox}

ホーム画面は 64行目で \texttt{useTime()} から値を受け取り、173行目の \texttt{\{todayMinutes\}} の場所に表示する。
ホーム画面のコードには「記録を読み込む」処理も「描き直す」処理も書いていない。
\textgt{Context の値が変わると、それを使っている画面は React が自動で描き直す}から（⑬）。
キャラクターの成長（累計 60 分・300 分・600 分・1500 分で次の姿になる）も、同じ \texttt{totalMinutes} から決まる。

\begin{kkbadgepoint}[要点]{この流れから分かる、アプリの基本の形}
  画面（\texttt{timer.tsx}・\texttt{index.tsx}）は、記録を自分では持たない。
  記録を変えたいときは Context の関数（\texttt{addMinutes} など）を呼び、表示したいときは Context の値（\texttt{todayMinutes} など）を受け取るだけ。
  保存は Context の中の \texttt{useEffect} が、値が変わるたびに自動で行う。
  だから「タイマー画面で足す」と「ホーム画面に出る」が、互いのコードを知らなくてもつながる。
  ほかの機能（疲労・設定）も、この同じ形でできている。
\end{kkbadgepoint}

\subsection{細部：出てきたライブラリと関数}

この章のコードに出てきたものだけをまとめる。「どこから来たか」で色を分けた。

\begin{longtblr}[entry=none,label=none]{
  colspec={Q[l,wd=3.3cm] Q[l,wd=2.4cm] X[l] Q[l,wd=3.4cm]},
  width=\linewidth,
  rowhead=1,
  row{1}={bg=kkdark, fg=white, font=\small\gtfamily},
  hlines={black!15},
  cells={font=\small},
}
名前 & どこから来たか & 何をするか & Pythonでいうと \\
\SetRow{bg=impAbg} \texttt{useState} & React & 覚えておく値と、それを変える関数を作る & インスタンス変数＋setter \\
\SetRow{bg=impAbg} \texttt{useEffect} & React & 指定した値が変わったときに処理を動かす & 値を見張るコールバック \\
\SetRow{bg=impAbg} \texttt{createContext}・\texttt{useContext} & React & 共有データの箱を作る・箱から取り出す & Flask の \texttt{g} \\
\SetRow{bg=impBbg} \texttt{useRef} & React & 描き直しても消えない入れ物 & インスタンス変数 \\
\SetRow{bg=impBbg} \texttt{useCallback} & React & 関数を毎回作り直さずに覚えておく & なし（速さのための工夫） \\
\SetRow{bg=impAbg} \texttt{AsyncStorage.getItem}・\texttt{setItem} & ライブラリ（async-storage） & 端末の保存箱から読む・書く & ファイルの読み書き \\
\SetRow{bg=impBbg} \texttt{JSON.stringify}・\texttt{JSON.parse} & JavaScript 標準 & データと文字列を相互に変換 & \texttt{json.dumps}・\texttt{json.loads} \\
\SetRow{bg=impBbg} \texttt{setInterval}・\texttt{clearInterval} & JavaScript 標準 & 一定時間ごとに処理を動かす・止める & \texttt{threading.Timer} \\
\SetRow{bg=impBbg} \texttt{Math.round} & JavaScript 標準 & 四捨五入 & \texttt{round}（ただし 2.5 は Python だと 2、JS だと 3） \\
\SetRow{bg=impBbg} \texttt{new Date().toISOString()} & JavaScript 標準 & 今の日時を世界標準時の文字列にする & \texttt{datetime.now(timezone.utc).isoformat()} \\
\SetRow{bg=impAbg} \texttt{todayKey}・\texttt{localDateKey} & 自作（\texttt{utils/date.ts}） & 端末の時刻で \texttt{"YYYY-MM-DD"} を作る & \texttt{date.today().isoformat()} \\
\SetRow{bg=impBbg} \texttt{genId} & 自作（\texttt{StudyLogContext.tsx}） & 時刻と乱数から記録の番号を作る & \texttt{uuid4()} の簡易版 \\
\SetRow{bg=impCbg} \texttt{View}・\texttt{Text}・\texttt{TouchableOpacity} & React Native & 箱・文字・押せる部品 & tkinter の \texttt{Frame}・\texttt{Label}・\texttt{Button} \\
\SetRow{bg=impCbg} \texttt{StyleSheet.create} & React Native & 見た目の設定をまとめる & tkinter の \texttt{configure()}（開発ログ V1） \\
\SetRow{bg=impCbg} \texttt{SafeAreaView} & ライブラリ（safe-area-context） & 画面の切り欠きを避けて表示する & なし \\
\SetRow{bg=impCbg} \texttt{=>}（アロー関数） & JavaScript の書き方 & 短く書ける関数 & \texttt{lambda} \\
\SetRow{bg=impCbg} \texttt{.{}.{}.}（スプレッド構文） & JavaScript の書き方 & 配列や辞書の中身を広げる & \texttt{*list}・\texttt{**dict} \\
\SetRow{bg=impCbg} \texttt{??} & JavaScript の書き方 & 左が空（\texttt{null}／\texttt{undefined}）なら右を使う & \texttt{x if x is not None else 0} \\
\SetRow{bg=impCbg} \texttt{条件 \&\& 部品} & JavaScript の書き方 & 条件が正しいときだけ部品を表示 & \texttt{if 条件: 表示} \\
\end{longtblr}


% =====================================================================
\section{例題}

\begin{kkbadgeexample}[四捨五入と、0 分の扱い]{問題1}
  タイマーで「開始」してから 25 秒で「停止」し、後チェックを完了した。
  ホームの「今日の学習時間」は何分増えるか。また、90 秒だった場合はどうか。
\end{kkbadgeexample}
\begin{kkbadgenote}[解答]
  \textgt{25 秒のとき：増えない（0 分）。}
  \texttt{timer.tsx} 86行目で \texttt{Math.round(25 / 60)}、つまり \texttt{Math.round(0.416…)} を計算して 0 になる。
  87行目で \texttt{addMinutes(0)} が呼ばれるが、\texttt{TimeContext.tsx} 66行目の \texttt{if (minutes <= 0) return;} で何もせずに終わる。
  ただし \texttt{StudySession} の記録は④で作られているので、\texttt{finishSession} によって \texttt{durationMinutes} が 0 の記録として残る。

  \textgt{90 秒のとき：2 分増える。}
  \texttt{Math.round(90 / 60)} は \texttt{Math.round(1.5)} で、JavaScript では 2 になる（Python の \texttt{round(1.5)} も 2 だが、\texttt{round(2.5)} は 2 なので注意）。
\end{kkbadgenote}

\begin{kkbadgeexample}[日付をまたぐとどうなるか]{問題2}
  23時50分に「開始」し、翌日の0時20分に「停止」した（ちょうど 30 分）。
  カレンダーでは、どちらの日に何分記録されるか。
\end{kkbadgeexample}
\begin{kkbadgenote}[解答]
  \textgt{停止した翌日に 30 分。}
  日付を決めているのは \texttt{TimeContext.tsx} 67行目の \texttt{const today = todayKey();} で、
  これは \texttt{addMinutes} が呼ばれたとき、つまり後チェックを完了した瞬間に実行される。
  そのため、開始した日ではなく停止した日に 30 分すべてが足され、前日には 1 分も入らない。

  なお、\texttt{todayKey()} は端末の時刻（日本時間）で日付を作る（\texttt{utils/date.ts}）。
  以前は世界標準時の日付を使っていたため、日本時間の 0〜9 時は前日として記録されていた。
  これを直したのがコミット \texttt{4b9acaa}。\evCode
\end{kkbadgenote}

\begin{kkbadgeexample}[アプリを閉じても消えない理由]{問題3}
  アプリを閉じて開き直しても、累計学習時間が 0 に戻らないのはなぜか。
  また、\texttt{TimeContext.tsx} の \texttt{loaded} がなかったら、何が起きる可能性があるか。
\end{kkbadgeexample}
\begin{kkbadgenote}[解答]
  \textgt{起動時に保存箱から読み戻しているから。}
  \texttt{TimeProvider} が最初に描かれたとき、次の \texttt{useEffect}（\texttt{[]} なので最初の1回だけ動く）が
  \texttt{timeLogData} を読み、\texttt{totalMinutes} と \texttt{history} に入れ直す。

\begin{codebox}{impA}{context/TimeContext.tsx}{36}
  useEffect(() => {
    (async () => {
      try {
        const raw = await AsyncStorage.getItem(STORAGE_KEY);
        if (raw) {
          const data = JSON.parse(raw);
          setTotalMinutes(data.totalMinutes ?? 0);

          if (data.history) {
            setHistory(data.history);
          } else if (data.savedDate && (data.todayMinutes ?? 0) > 0) {
            // 旧フォーマットからの移行
            setHistory({ [data.savedDate]: data.todayMinutes });
          }
        }
      } catch {
        // 読み込み失敗時はデフォルト値のまま
      } finally {
        setLoaded(true);
      }
    })();
  }, []);
\end{codebox}

  \textgt{\texttt{loaded} がないと、保存していた記録が 0 で上書きされるおそれがある。}
  起動した直後は \texttt{totalMinutes} が 0、\texttt{history} が空の状態で始まる。
  保存の \texttt{useEffect}（59〜63行目）は値が変わるたびに動くので、読み込みが終わる前にこの「空の状態」を保存してしまうかもしれない。
  60行目の \texttt{if (!loaded) return;} は、54行目で読み込みが終わったと分かるまで保存を止めておくためのもの。
  この理由は開発ログ V1 の Step 5 にも書かれている。\evLog
\end{kkbadgenote}

% =====================================================================
\section{まとめ}

\subsection{今回の結論}
\begin{itemize}
  \item このアプリは TypeScript で書かれ、React Native と Expo によって iPhone・Android・Web の3つで動く。記録は端末の中（AsyncStorage）にだけ保存し、Supabase はログインだけに使う。
  \item 独自の中心は \texttt{app/(tabs)/}（画面）・\texttt{context/}（共有データと保存）・\texttt{types/}（データの形）の3か所。
  \item 画面は記録を持たず、Context の関数を呼ぶだけ。保存と描き直しは自動で起きる。
\end{itemize}

\subsection{復習するところ}
\begin{itemize}
  \item \tagA\ 図5の①〜⑬を、コードを見ずに言えるか。
  \item \tagA\ \texttt{context/TimeContext.tsx}（90行）を最初から最後まで読む。今回ほぼ全部を見ている。
  \item \tagB\ \texttt{types/studyLog.ts} と図4を見比べ、項目の意味を言えるか。
  \item \tagB\ 定義の章の React の用語（state・\texttt{useEffect}・Context）を、Python の言葉で説明できるか。
\end{itemize}

\subsection{次に持ち越す疑問}
\begin{enumerate}
  \item \textgt{疲労スコアはどう計算されているか。}\texttt{utils/fatigueCalc.ts} と \texttt{hooks/use-fatigue.ts} は今回読んでいない。同じ「縦に追う」やり方で読める。
  \item \textgt{学習時間が2か所に記録されているのは意図どおりか。}\texttt{timeLogData}（累計と日ごと）と、\texttt{studyLogData\_v1} の各 \texttt{StudySession} の \texttt{durationMinutes} の両方に同じ分数が入る。
  \item \textgt{機種を変えたら記録はどうなるか。}記録は端末の中にしかないので、今のままでは引き継げない。
  \item \textgt{アプリを裏に回している間も、タイマーは正しく数えるか。}1秒ごとに \texttt{seconds} を足す仕組みなので、スマホがアプリを止めている間は数えられない可能性がある（未確認）。
\end{enumerate}



\end{document}
